\documentclass[10pt,letterpaper]{article}
\usepackage[utf8]{inputenc}
\usepackage[intlimits]{amsmath}
\usepackage{amsfonts}
\usepackage{amssymb}
\usepackage{ragged2e}
\usepackage[letterpaper, margin=0.5in]{geometry}
\usepackage{graphicx}
\usepackage{cancel}
\usepackage{mathtools}
\usepackage{tabularx}
\usepackage{arydshln}
\usepackage{tensor}
\usepackage{array}
\usepackage{xcolor}
\usepackage[boxed]{algorithm}
\usepackage[noend]{algpseudocode}
\usepackage{listings}
\usepackage{textcomp}
% \usepackage[pdf,tmpdir,singlefile]{graphviz}
\usepackage{mathrsfs}
\usepackage{bbm}
\usepackage{tikz}
\usepackage{tikz-cd}
\usepackage{enumitem}
\usepackage{arydshln}
\usepackage{relsize}
\usepackage{multicol}
\usepackage{scalerel}
\usepackage{upgreek}
\usepackage{wasysym}

\usetikzlibrary{bayesnet}
\setlist{noitemsep}

%%%%%%%%%%%%%%%%%%%%%%%%%%%%%
% Formatting commands
%%%%%%%%%%%%%%%%%%%%%%%%%%%%%
\newcommand{\n}{\hfill\break}
\newcommand{\nn}{\vspace{0.5\baselineskip}\n}
\newcommand{\up}{\vspace{-\baselineskip}}
\newcommand{\hangblock}[2]{\par\noindent\settowidth{\hangindent}{\textbf{#1: }}\textbf{#1: }\nolinebreak#2}
\newcommand{\lemma}[1]{\hangblock{Lemma}{#1}}
\newcommand{\defn}[1]{\hangblock{Defn}{#1}}
\newcommand{\thm}[1]{\hangblock{Thm}{#1}}
\newcommand{\cor}[1]{\hangblock{Cor}{#1}}
\newcommand{\prop}[1]{\hangblock{Prop}{#1}}
\newcommand{\ex}[1]{\hangblock{Ex}{#1}}
\newcommand{\exer}[1]{\hangblock{Exer}{#1}}
\newcommand{\fact}[1]{\hangblock{Fact}{#1}}
\newcommand{\remark}[1]{\hangblock{Remark}{#1}}
\newcommand{\proven}{\;$\square$\n}
\newcommand{\problem}[1]{\par\noindent{#1}\n}
\newcommand{\problempart}[2]{\par\noindent\indent{}\settowidth{\hangindent}{\textbf{(#1)} \indent{}}\textbf{(#1) }\nolinebreak#2\n}
\newcommand{\ptxt}[1]{\textrm{\textnormal{#1}}}
\newcommand{\inlineeq}[1]{\centerline{$\displaystyle #1$}}
\newcommand{\pageline}{\up\par\noindent\rule{\textwidth}{0.1pt}}

%%%%%%%%%%%%%%%%%%%%%%%%%%%%%
% Math commands
%%%%%%%%%%%%%%%%%%%%%%%%%%%%%
% Set Theory
\newcommand{\card}[1]{\left|#1\right|}
\newcommand{\set}[1]{\left\{#1\right\}}
\newcommand{\setmid}{\;\middle|\;}
\newcommand{\ps}[1]{\mathcal{P}\left(#1\right)}
\newcommand{\pfinite}[1]{\mathcal{P}^{\ptxt{finite}}\left(#1\right)}
\newcommand{\naturals}{\mathbb{N}}
\newcommand{\N}{\naturals}
\newcommand{\integers}{\mathbb{Z}}
\newcommand{\Z}{\integers}
\newcommand{\rationals}{\mathbb{Q}}
\newcommand{\Q}{\rationals}
\newcommand{\reals}{\mathbb{R}}
\newcommand{\R}{\reals}
\newcommand{\complex}{\mathbb{C}}
\newcommand{\C}{\complex}
\newcommand{\halfPlane}{\mathbb{H}}
\let\HSym\H
\let\H\relax
\newcommand{\H}{\halfPlane}
\newcommand{\comp}{^{\complement}}
\DeclareMathOperator{\Hom}{Hom}
\newcommand{\Ind}{\mathbbm{1}}
\newcommand{\cut}{\setminus}
\DeclareMathOperator{\elem}{elem}

% Differentiable Manifolds
\newcommand{\RP}{\mathbb{RP}}
\newcommand{\CP}{\mathbb{RP}}
\newcommand{\osubset}{\overset{\mathclap{\scalebox{0.5}{\ptxt{open}}}}{\subset}}
\newcommand{\osubseteq}{\overset{\mathclap{\scalebox{0.5}{\ptxt{open}}}}{\subseteq}}
\newcommand{\osupset}{\overset{\mathclap{\scalebox{0.5}{\ptxt{open}}}}{\supset}}
\newcommand{\osupseteq}{\overset{\mathclap{\scalebox{0.5}{\ptxt{open}}}}{\supseteq}}
\newcommand{\pdat}[3]{\left.\pd{#1}{#2}\right|_{#3}}
\DeclareMathOperator{\so}{so}
\DeclareMathOperator{\codim}{codim}
\DeclareMathOperator{\Diff}{Diff}
\let\dSym\d
\let\d\relax
\newcommand{\d}{\partial}
\DeclareMathOperator{\gl}{gl}

% Graph Theory
\let\deg\relax
\DeclareMathOperator{\deg}{deg}
\newcommand{\degp}{\ptxt{deg}^{+}}
\newcommand{\degn}{\ptxt{deg}^{-}}
\newcommand{\precdot}{\mathrel{\ooalign{$\prec$\cr\hidewidth\hbox{$\cdot\mkern0.5mu$}\cr}}}
\newcommand{\succdot}{\mathrel{\ooalign{$\cdot\mkern0.5mu$\cr\hidewidth\hbox{$\succ$}\cr\phantom{$\succ$}}}}
\DeclareMathOperator{\cl}{cl}
\DeclareMathOperator{\affdim}{affdim}

% Probability
\newcommand{\parSymbol}{\P}
\newcommand{\Prob}{\mathbb{P}}
\renewcommand{\P}{\Prob}
\newcommand{\Avg}{\mathbb{E}}
\newcommand{\E}{\Avg}
\DeclareMathOperator{\Var}{Var}
\DeclareMathOperator{\cov}{cov}
\DeclareMathOperator{\Unif}{Unif}
\DeclareMathOperator{\Binom}{Binom}
\newcommand{\CI}{\mathrel{\text{\scalebox{1.07}{$\perp\mkern-10mu\perp$}}}}

% Standard Math
\newcommand{\inv}{^{-1}}
\newcommand{\abs}[1]{\left|#1\right|}
\newcommand{\ceil}[1]{\left\lceil{}#1\right\rceil{}}
\newcommand{\floor}[1]{\left\lfloor{}#1\right\rfloor{}}
\newcommand{\conj}[1]{\overline{#1}}
\newcommand{\of}{\circ}
\newcommand{\tri}{\triangle}
\newcommand{\inj}{\hookrightarrow}
\newcommand{\surj}{\twoheadrightarrow}
\newcommand{\ndiv}{\nmid}
\renewcommand{\epsilon}{\varepsilon}
\newcommand{\divides}{\mid}
\newcommand{\ndivides}{\nmid}
\DeclareMathOperator{\lcm}{lcm}
\DeclareMathOperator{\sgn}{sgn}
\newcommand{\map}[4]{\!\!\!\begin{array}[t]{rcl}#1 & \!\!\!\!\to & \!\!\!\!#2\\ {}#3 & \!\!\!\!\mapsto & \!\!\!\!#4\end{array}}
\newcommand{\bigsum}[2]{\smashoperator[lr]{\sum_{\scalebox{#1}{$#2$}}}}
\DeclareMathOperator{\gcf}{gcf}
\newcommand{\restr}[2]{\left.#1\right|_{#2}}

% Linear Algebra
\newcommand{\Id}{\textrm{\textnormal{Id}}}
\newcommand{\im}{\textrm{\textnormal{im}}}
\newcommand{\norm}[1]{\abs{\abs{#1}}}
\newcommand{\tpose}{^{T}\!}
\newcommand{\iprod}[1]{\left<#1\right>}
\newcommand{\giprod}{\iprod{\;\,,\;}}
\DeclareMathOperator{\tr}{tr}
\DeclareMathOperator{\trace}{tr}
\newcommand{\chgBasMat}[3]{\!\!\tensor*[_{#1}]{\left[#2\right]}{_{#3}}}
\newcommand{\vecBas}[2]{\tensor*[]{\left[#1\right]}{_{#2}}}
\DeclareMathOperator{\GL}{GL}
\DeclareMathOperator{\Mat}{Mat}
\DeclareMathOperator{\vspan}{span}
\DeclareMathOperator{\rank}{rank}
\newcommand{\V}[1]{\vec{#1}}
\DeclareMathOperator{\proj}{proj}
\DeclareMathOperator{\compProj}{comp}
\DeclareMathOperator{\row}{row}
\newcommand{\smallPMatrix}[1]{\paren{\begin{smallmatrix}#1\end{smallmatrix}}}
\newcommand{\smallBMatrix}[1]{\brack{\begin{smallmatrix}#1\end{smallmatrix}}}
\newcommand{\pmat}[1]{\begin{pmatrix}#1\end{pmatrix}}
\newcommand{\bmat}[1]{\begin{bmatrix}#1\end{bmatrix}}
\newcommand{\dual}{^{*}}
\newcommand{\pinv}{^{\dagger}}
\newcommand{\horizontalMatrixLine}{\ptxt{\rotatebox[origin=c]{-90}{$|$}}}
\DeclareMathOperator{\range}{range}
\DeclareMathOperator{\Symm}{Symm}
\DeclareMathOperator{\SU}{SU}

% Multilinear Algebra
\let\Lsym\L
\let\L\relax
\DeclareMathOperator{\L}{\mathscr{L}}
\DeclareMathOperator{\A}{\mathcal{A}}
\DeclareMathOperator{\Alt}{Alt}
\DeclareMathOperator{\Sym}{Sym}
\newcommand{\ot}{\otimes}
\newcommand{\ox}{\otimes}
\DeclareMathOperator{\asc}{asc}
\DeclareMathOperator{\asSet}{set}
\DeclareMathOperator{\sort}{sort}
\DeclareMathOperator{\ringA}{\mathring{A}}
\DeclareMathOperator{\Sh}{Sh}

% Topology
\newcommand{\closure}[1]{\overline{#1}}
\newcommand{\uball}{\mathcal{U}}
\DeclareMathOperator{\Int}{Int}
\DeclareMathOperator{\Ext}{Ext}
\DeclareMathOperator{\Bd}{Bd}
\DeclareMathOperator{\rInt}{rInt}
\DeclareMathOperator{\ch}{ch}
\DeclareMathOperator{\ah}{ah}
\newcommand{\LargerTau}{\mathlarger{\mathlarger{\mathlarger{\mathlarger{\tau}}}}}
\newcommand{\Tau}{\mathcal{T}}

% Analysis
\DeclareMathOperator{\Graph}{Graph}
\DeclareMathOperator{\epi}{epi}
\DeclareMathOperator{\hypo}{hypo}
\DeclareMathOperator{\supp}{supp}
\newcommand{\lint}[2]{\underset{#1}{\overset{#2}{{\color{black}\underline{{\color{white}\overline{{\color{black}\int}}\color{black}}}}}}}
\newcommand{\uint}[2]{\underset{#1}{\overset{#2}{{\color{white}\underline{{\color{black}\overline{{\color{black}\int}}\color{black}}}}}}}
\newcommand{\alignint}[2]{\underset{#1}{\overset{#2}{{\color{white}\underline{{\color{white}\overline{{\color{black}\int}}\color{black}}}}}}}
\newcommand{\extint}{\ptxt{ext}\int}
\newcommand{\extalignint}[2]{\ptxt{ext}\alignint{#1}{#2}}
\newcommand{\conv}{\ast}
\newcommand{\pd}[2]{\frac{\partial{}#1}{\partial{}#2}}
\newcommand{\del}{\nabla}
\DeclareMathOperator{\grad}{grad}
\DeclareMathOperator{\curl}{curl}
\let\div\relax
\DeclareMathOperator{\div}{div}
\DeclareMathOperator{\vol}{vol}

% Complex Analysis
\let\Re\relax
\DeclareMathOperator{\Re}{Re}
\let\Im\relax
\DeclareMathOperator{\Im}{Im}
\DeclareMathOperator{\Res}{Res}

% Abstract Algebra
\DeclareMathOperator{\ord}{ord}
\newcommand{\generated}[1]{\left<#1\right>}
\newcommand{\cycle}[1]{\smallPMatrix{#1}}
\newcommand{\id}{\ptxt{id}}
\newcommand{\iso}{\cong}
\DeclareMathOperator{\Aut}{Aut}
\DeclareMathOperator{\SL}{SL}
\DeclareMathOperator{\op}{op}
\newcommand{\isom}[4]{\!\!\!\begin{array}[t]{rcl}#1 & \!\!\!\!\overset{\sim}{\to} & \!\!\!\!#2\\ #3 & \!\!\!\!\mapsto & \!\!\!\!#4\end{array}}
\newcommand{\F}{\mathbb{F}}
\newcommand{\acts}{\;\reflectbox{\rotatebox[origin=c]{-90}{$\circlearrowright$}}\;}
\newcommand{\disjunion}{\mathrel{\text{\scalebox{1.07}{$\perp\mkern-10mu\perp$}}}}
\DeclareMathOperator{\SO}{SO}
\DeclareMathOperator{\stab}{stab}
\DeclareMathOperator{\nullity}{nullity}
\DeclareMathOperator{\Perm}{Perm}
\DeclareMathOperator{\nsubgp}{\vartriangleleft}
\DeclareMathOperator{\notnsubgp}{\ntriangleleft}
\newcommand{\presentation}[2]{\left<#1\;\middle|\;#2\right>}
\DeclareMathOperator{\Char}{char}
\DeclareMathOperator{\fchar}{char}
\DeclareMathOperator{\triv}{triv}
\DeclareMathOperator{\reg}{reg}
\DeclareMathOperator{\std}{std}
\DeclareMathOperator{\Func}{Func}
\DeclareMathOperator{\End}{End}

% Convex Optimization
\let\sectionSymbol\S
\let\S\relax
\newcommand{\S}{\mathbb{S}}
\DeclareMathOperator{\dist}{dist}
\DeclareMathOperator{\dom}{dom}
\DeclareMathOperator{\diag}{diag}
\DeclareMathOperator{\ones}{\mathbbm{1}}
\newcommand{\minimizeOver}[3]{\begin{array}{rl}\underset{#1}{\ptxt{minimize}} & #2\\ \ptxt{subject to} & #3\end{array}}
\newcommand{\maximizeOver}[3]{\begin{array}{rl}\underset{#1}{\ptxt{maximize}} & #2\\ \ptxt{subject to} & #3\end{array}}
\newcommand{\minimizationProblem}[2]{\minimizeOver{}{#1}{#2}}
\newcommand{\maximizationProblem}[2]{\maximizeOver{}{#1}{#2}}
\newcommand{\minimizeOverUnconstrained}[2]{\begin{array}{rl}\underset{#1}{\ptxt{minimize}} & #2\end{array}}
\newcommand{\maximizeOverUnconstrained}[2]{\begin{array}{rl}\underset{#1}{\ptxt{maximize}} & #2\end{array}}
\newcommand{\minimizationUnconstrained}[1]{\minimizeOverUnconstrained{}{#1}}
\newcommand{\maximizationUnconstrained}[1]{\maximizeOverUnconstrained{}{#1}}
\DeclareMathOperator{\argmin}{argmin}
\DeclareMathOperator{\argmax}{argmax}

% Proofs
\newcommand{\st}{s.t.}
\newcommand{\unique}{!}
\newcommand{\iffdef}{\overset{\ptxt{def}}{\Leftrightarrow}}
\newcommand{\eqVertical}{\rotatebox[origin=c]{90}{=}}
\newcommand{\mapsfrom}{\mathrel{\reflectbox{\ensuremath{\mapsto}}}}
\newcommand{\mapsdown}{\rotatebox[origin=c]{-90}{$\mapsto$}\mkern2mu}
\newcommand{\mapsup}{\rotatebox[origin=c]{90}{$\mapsto$}\mkern2mu}
\newcommand{\from}{\!\mathrel{\reflectbox{\ensuremath{\to}}}}
\newcommand{\labeledeq}[1]{\overset{\mathclap{\ptxt{#1}}}{=}}
\newcommand{\eqdef}{\labeledeq{def}}

% Brackets
\newcommand{\paren}[1]{\left(#1\right)}
\renewcommand{\brack}[1]{\left[#1\right]}
\renewcommand{\brace}[1]{\left\{#1\right\}}
\newcommand{\ang}[1]{\left<#1\right>}

% Algorithms
\algrenewcommand{\algorithmiccomment}[1]{\hskip 1em \texttt{// #1}}
\algrenewcommand\algorithmicrequire{\textbf{Input:}}
\algrenewcommand\algorithmicensure{\textbf{Output:}}
\newcommand{\algP}{\ptxt{\textbf{P}}}
\newcommand{\algNP}{\ptxt{\textbf{NP}}}
\newcommand{\algNPC}{\ptxt{\textbf{NP-Complete}}}
\newcommand{\algNPH}{\ptxt{\textbf{NP-Hard}}}
\newcommand{\algEXP}{\ptxt{\textbf{EXP}}}
\DeclareMathOperator{\fl}{fl}

%%%%%%%%%%%%%%%%%%%%%%%%%%%%%
% Other commands
%%%%%%%%%%%%%%%%%%%%%%%%%%%%%
\newcommand{\flag}[1]{\textbf{\textcolor{red}{#1}}}
\let\uSym\u
\let\u\relax
\newcommand{\u}[1]{\underline{#1}}
\let\bSym\b
\let\b\relax
\newcommand{\b}[1]{\textbf{#1}}
\let\iSym\i
\let\i\relax
\newcommand{\i}[1]{\textit{#1}}
\let\scSym\sc
\let\sc\relax
\newcommand{\sc}[1]{\textsc{#1}}

%%%%%%%%%%%%%%%%%%%%%%%%%%%%%%%%%%%%%%%
% Make l's curvy in math environments %
%%%%%%%%%%%%%%%%%%%%%%%%%%%%%%%%%%%%%%%
\mathcode`l="8000
\begingroup
\makeatletter
\lccode`\~=`\l
\DeclareMathSymbol{\lsb@l}{\mathalpha}{letters}{`l}
\lowercase{\gdef~{\ifnum\the\mathgroup=\m@ne \ell \else \lsb@l \fi}}%
\endgroup

%%%%%%%%%%%%%%%%%%%%%%%%%
% Fix \vdots and \ddots %
%%%%%%%%%%%%%%%%%%%%%%%%%
\usepackage{letltxmacro}
\LetLtxMacro\orgvdots\vdots
\LetLtxMacro\orgddots\ddots

\makeatletter
\DeclareRobustCommand\vdots{%
	\mathpalette\@vdots{}%
}
\newcommand*{\@vdots}[2]{%
	% #1: math style
	% #2: unused
	\sbox0{$#1\cdotp\cdotp\cdotp\m@th$}%
	\sbox2{$#1.\m@th$}%
	\vbox{%
		\dimen@=\wd0 %
		\advance\dimen@ -3\ht2 %
		\kern.5\dimen@
		% remove side bearings
		\dimen@=\wd2 %
		\advance\dimen@ -\ht2 %
		\dimen2=\wd0 %
		\advance\dimen2 -\dimen@
		\vbox to \dimen2{%
			\offinterlineskip
			\copy2 \vfill\copy2 \vfill\copy2 %
		}%
	}%
}
\DeclareRobustCommand\ddots{%
	\mathinner{%
		\mathpalette\@ddots{}%
		\mkern\thinmuskip
	}%
}
\newcommand*{\@ddots}[2]{%
	% #1: math style
	% #2: unused
	\sbox0{$#1\cdotp\cdotp\cdotp\m@th$}%
	\sbox2{$#1.\m@th$}%
	\vbox{%
		\dimen@=\wd0 %
		\advance\dimen@ -3\ht2 %
		\kern.5\dimen@
		% remove side bearings
		\dimen@=\wd2 %
		\advance\dimen@ -\ht2 %
		\dimen2=\wd0 %
		\advance\dimen2 -\dimen@
		\vbox to \dimen2{%
			\offinterlineskip
			\hbox{$#1\mathpunct{.}\m@th$}%
			\vfill
			\hbox{$#1\mathpunct{\kern\wd2}\mathpunct{.}\m@th$}%
			\vfill
			\hbox{$#1\mathpunct{\kern\wd2}\mathpunct{\kern\wd2}\mathpunct{.}\m@th$}%
		}%
	}%
}
\makeatother

\newcommand{\B}{
	\begin{tikzpicture}
	\filldraw [fill=red, draw=black] (0, 0) rectangle (0.37, 0.45);
	\draw [line width=0.5mm, white ] (0.1,0.08) -- (0.1,0.38);
	\draw[line width=0.5mm, white ] (0.1, 0.35) .. controls (0.2, 0.35) and (0.4, 0.2625) .. (0.1, 0.225);
	\draw[line width=0.5mm, white ] (0.1, 0.225) .. controls (0.2, 0.225) and (0.4, 0.1625) .. (0.1, 0.1);
	\end{tikzpicture}
}

% Allow custom symbols as arrows in tikz-cd
\tikzset{
  symbol/.style={
    draw=none,
    every to/.append style={
      edge node={node [sloped, allow upside down, auto=false]{$#1$}}}
  }
}
% Example: \arrow[r,symbol=\cong]
% https://tex.stackexchange.com/questions/394154/how-to-include-inclusion-subgroup-relationship-in-tikz-cd-diagram

\author{Thomas Cohn}
\title{Math 591 Lecture 33}
\date{11/18/20} % Can also use \today

\begin{document}
\maketitle
\setlength\RaggedRightParindent{\parindent}
\RaggedRight

\section*{Global formulae for $d$}

\par\noindent
First, we want to consider the pairing of forms and vector fields. If $\alpha\in\Omega^{k}$ and $X_{1},\ldots,X_{k}\in\mathfrak{X}(M)$, then $\alpha(X_{1},\ldots,X_{k}):M\to\R$ is a function on $M$.\n

\prop{
	(1) If $\alpha\in\Omega^{1}(M)$, $Y_{0},Y_{1}\in\mathfrak{X}(M)$, then
	\[
		(d\alpha)(Y_{0},Y_{1})=Y_{0}(\alpha(Y_{1}))-Y_{1}(\alpha(Y_{0}))-\alpha(\brack{Y_{0},Y_{1}})
	\]
	Proof: WOLOG assume $\alpha=f\,dg$. Then $d\alpha=df\wedge{}dg$. Thus,
	\[
		(d\alpha)(Y_{0},Y_{1})=df(Y_{0})dg(Y_{1})-df(Y_{1})dg(Y_{0})=Y_{0}(f)Y_{1}(g)-Y_{1}(f)Y_{0}(g)
	\]
	Well, the RHS of this is
	\[
		Y_{0}(fY_{1}(g))-Y_{1}(fY_{0}(g))-f\brack{Y_{0},Y_{1}}g=(Y_{0}f)(Y_{1}g)+\cancel{fY_{0}Y_{1}g}-(Y_{1}fY_{0}g+\cancel{fY_{1}Y_{0}g})-f(\cancel{Y_{0}Y_{1}g}-\cancel{Y_{1}Y_{0}g})
	\]
	\proven
}

\prop{
	(2) If $\alpha\in\Omega^{2}(M)$, $Y_{0},Y_{1},Y_{2}\in\mathfrak{X}(M)$, then
	\[
		(d\alpha)(Y_{0},Y_{1},Y_{2})=Y_{0}\alpha(Y_{1},Y_{2})-Y_{1}\alpha(Y_{0},Y_{2})+Y_{2}\alpha(Y_{0},Y_{1})-\alpha(\brack{Y_{0},Y_{1}},Y_{2})+\alpha(\brack{Y_{0},Y_{2}},Y_{1})-\alpha(\brack{Y_{1},Y_{2}},Y_{0})
	\]
	In general, if $\alpha\in\Omega^{k}(M)$, $Y_{0},\ldots,Y_{k}\in\mathfrak{X}(M)$, then
	\[
		(d\alpha)(Y_{0},\ldots,Y_{k})=\sum_{j=0}^{k}(-1)^{j}Y_{j}\alpha(Y_{0},\ldots,\widehat{Y_{j}},\ldots,Y_{k})+\smashoperator{\sum_{0\le{}i<j\le{}k}}(-1)^{i+j}\alpha(\brack{Y_{i},Y_{j}},Y_{0},\ldots,\widehat{Y_{i}},\ldots,\widehat{Y_{j}},\ldots,Y_{k})
	\]
	where $\widehat{\;\,\cdot\;\,}$ indicates that we omit that term.\n
}

\par\noindent
Refer to Proposition 20.14 in Tu/Proposition 14.32 in Lee.\n

\par\noindent
Another way to think of this is, gven $\alpha\in\Omega^{k}(M)$, $X_{1},\ldots,X_{k}\in\mathfrak{X}(M)$, and $p\in{}M$, we have $(X_{i})_{p}\in{}T_{p}M$, so
\[
	\map{\alpha_{p}:\smash{\underbrace{T_{p}M\times\cdots\times{}T_{p}M}_{k}}=(T_{p}M)^{k}}{\R}{p}{\alpha_{p}((X_{1})_{p},\ldots,(X_{k})_{p})\in\R}
\]

\section*{Lie Derivatives of Forms}

\defn{
	Let $X\in\mathfrak{X}(M)$, $\alpha\in\Omega^{k}(M)$, and $p\in{}M$. Let $\varphi$ be the local flow of $X$. Form the curve $(-\varepsilon,\varepsilon)\ni{}t\mapsto(\varphi_{t}^{*}\alpha)_{p}\in\bigwedge^{k}(T_{p}^{*}M)$. (Note that $\bigwedge^{k}(T_{p}^{*}M)$ is independent of $t$.) Then we defined
	\[
		(\mathcal{L}_{X}\alpha)_{p}=\restr{\frac{d}{dt}(\varphi_{t}^{*}\alpha)_{p}}{t=0}
	\]
}

\par\noindent
Observe: If $k=0$, then $\alpha$ is a function on $M$, so
\[
	(\mathcal{L}_{X}\alpha)_{p}=\restr{\frac{d}{dt}(\alpha\of\varphi_{t})_{p}}{t=0}=\restr{\frac{d}{dt}(\alpha(\varphi_{t}(p)))}{t=0}=\restr{\dot\varphi_{t}\vphantom{\frac{d}{dt}}}{\mathrlap{t=0}}(\alpha)
\]
Thus, $(\mathcal{L}_{X}\alpha)_{p}=X(\alpha)(p)$ if $\alpha$ is a function.\n

\subsection*{First Properties}

\prop{
	\begin{enumerate}[leftmargin=4\parindent, topsep=0pt, itemsep=0pt, label=\alph*)]
		\item The expression $\mathcal{L}_{X}(\alpha)$ is $\R$-linear w.r.t $X$ and $\alpha$.
		\item $\mathcal{L}_{X}(\alpha\wedge\beta)=\mathcal{L}_{X}(\alpha)\wedge\beta+\alpha\wedge\mathcal{L}_{X}(\beta)$ -- $\mathcal{L}_{X}$ is a derivation.
		\item $\mathcal{L}_{X}(d\alpha)=d\mathcal{L}_{X}(\alpha)$.
		\item If $F:M\to{}M$ is a diffeomorphism, then $\mathcal{L}_{X}(F^{*}\alpha)=F^{*}\paren{\mathcal{L}_{F_{*}(X)}\alpha}$.
	\end{enumerate}\up\nn
	Proof: (b) holds because $\varphi_{t}^{*}(\alpha\wedge\beta)=(\varphi_{t}^{*}\alpha)\wedge(\varphi_{t}^{*}\beta)$. We just need to show the right hand side satisfies he usual product rule with respect to $\frac{d}{dt}$.\n
	(c) holds because $\varphi_{t}^{*}(d\alpha)=d(\varphi_{t}^{*}\alpha)$, $\forall{}t$. This uses the fact that $d$ and $\frac{d}{dt}$ commute, i.e., $\frac{\d^{2}}{\d{}t\d{}x^{j}}=\frac{\d^{2}}{\d{}x^{j}\d{}t}$.\n
	(d) holds because for $p\in{}M$, $v_{1},\ldots,v_{k}\in{}T_{p}M$,
	\[
		\mathcal{L}_{X}(F^{*}\alpha)_{p}(v_{1},\ldots,v_{k})=\restr{\frac{d}{dt}\smash{\underbrace{\varphi_{t}^{*}(F^{*}\alpha)_{p}\vphantom{\frac{d}{dt}}}_{\mathclap{(\varphi_{t}^{*}\of{}F^{*})(\alpha)=(F\of\varphi_{t})^{*}(\alpha)}}}(v_{1},\ldots,v_{k})}{t=0}=\restr{\frac{d}{dt}\alpha((F\of\varphi_{t})_{*,p})(v_{1}),\ldots,(F\of\varphi_{t})_{*,p}(v_{k})}{t=0}
	\]
	Observe: $t\mapsto{}F\of\varphi_{t}$ is the integral curve of $F_{*}(X)$ at $F(p)$. So
	\[
		\mathcal{L}_{X}(F^{*}\alpha)_{p}(v_{1},\ldots,v_{k})=\cdots=(\mathcal{L}_{F_{*}X}\alpha)(F_{*,p}(v_{1}),\ldots,F_{*,p}(v_{k}))=F^{*}(\mathcal{L}_{F_{*}(X)}\alpha)(v_{1},\ldots,v_{k})
	\]
	\proven
}

\cor{
	If $\mathcal{L}_{X}\alpha=0$ (everywhere), then $\varphi_{t}^{*}\alpha=\alpha$, $\forall{}t$.\nn
	Proof: We will show $\forall{}p\in{}M$, the curve $t\mapsto{}(\varphi_{t}^{*}\alpha)_{p}\in\bigwedge^{k}T_{p}^{*}M$ is constant.\nn
	Proof: $\forall{}s$,
	\[
		\restr{\frac{d}{dt}\varphi_{t}^{*}\alpha}{t=s}=\restr{\frac{d}{dt}(\varphi_{t+s}^{*}\alpha)}{t=0}=\restr{\frac{d}{dt}\varphi_{t}^{*}(\varphi_{s}^{*}(\alpha))}{t=0}=\mathcal{L}_{X}(\varphi_{s}^{*}\alpha)=\varphi_{s}^{*}\mathcal{L}_{\underbrace{(\varphi_{s})_{*}(X)}_{=X}}\alpha=\varphi_{s}^{*}\underbrace{\mathcal{L}_{X}\alpha}_{=0}=0
	\]
	\proven
}

\subsection*{Towards Cartan's Formula}

\par\noindent
We need interior multiplication/contraction.\n

\defn{
	Given $\alpha\in\Omega^{k}$, $X\in\mathfrak{X}(M)$, we define $\iota_{X}\alpha\in\Omega^{k-1}$ by, $\forall{}p\in{}M$ and $v_{1},\ldots,v_{k-1}\in{}T_{p}M$,
	\[
		(\iota_{X}\alpha)_{p}(v_{1},\ldots,v_{k-1})=\alpha_{p}(X_{p},v_{1},\ldots,v_{k})
	\]
}

\thm{
	(Cartan's Magic Formula) $\mathcal{L}_{X}=d\of\iota_{X}+\iota_{X}\of{}d:\Omega^{k}\to\Omega^{k}$.\n
}

\par\noindent
Picture:
\[
	\begin{tikzcd}
		\cdots \arrow[r,"d"] & \Omega^{k-1} \arrow[r,"d"] & \Omega^{k} \arrow[r,"d"] \arrow[dl,"\iota_{X}",swap] \arrow[d,"\mathcal{L}_{X}", Rightarrow] & \Omega^{k+1} \arrow[r,"d"] \arrow[dl,"\iota_{X}"] & \cdots\\
		\cdots \arrow[r,"d"] & \Omega^{k-1} \arrow[r,"d"] & \Omega^{k} \arrow[r,"d"] & \Omega^{k+1} \arrow[r,"d"] & \cdots
	\end{tikzcd}
\]
where $\mathcal{L}_{X}$ is the ``sum of two paths''.\n

\par\noindent
On Friday, we'll show that when $\alpha=dx\wedge{}dy$ and $M=\R^{2}$, then $\mathcal{L}_{X}(\alpha)=(\div{}X)\alpha$, in the calc III sense.

\end{document}